\subsection{Motivation}

Today more than ever, Recommendation systems became suggestion systems, losing the title of item sorter. They now influence the users based on their personal interest, on social media, ecommerce websites or movies to watch for example. When a Large Language model is added into the mix, the effect is more impactful. Recent studies show that models can be used to express the intent of a user, generate responses based on them, assist in customer service scenarios and also to rerank results \cite{wu2024llmrecsurvey,zhao2024llmrecsystems}. When using LLM-powered ranking, machines can make decisions alone based on old user interactions, item metadata without ever being supervised \cite{hou2024llmrankers}. This approach simplifies the recommendation process since theres no need for actual ranking criteria. Retrieval is now the new workflow, where it pulls the relevant items befoore sharing them with the LLM. The system is simple and logical, and solves a massive issue for corporations.

retrieval is a core part of infrastructure since its the layer that shapes the data received by the LLM. the external made chunks are what shapes the models output creating a clear path for Tampering. Instead of directly affecting the model, an attacker can corrupt the stored data, shift the logic or even modify the data before it ever gets to the model. This type of attack bypasses the traditional safeguards entirely, leaving room for poisoning of recommendation systems. Studies confirm that poisoned data or even modified metadata can completely reshape the response. The issue doesn't reside in the code itself but rather the data coming from the retrieval process as even a subtle change can greatly impact the recommendation quality through hidden instructions in payloads meant to disrupt the initial behavior of the pipeline. studies about prompt injection discuss this issue in detail \cite{greshake2023indirectpi,yi2025bipia}. This type of vulnerability is especially relevant in recommender systems since their functionning depends on the external data they retrieve.

This is why our study will discuss the engineering behind a red teaming through the trade offs that come from using LLMs in a recommendation system. RAGPoison Lab will develop on this idea using the very famous MovieLens 100k dataset to simulate the require environnement locally. We will compare the recommendation offered from a clean index to the one coming from a poisoned source. Every step of the workflow is logged, organized and audited to give the best analysis possible of the situation mesurin impact and interpreting the results in the hope of discussing defense options.

\subsection{Problem Statement}

This thesis doesn't just rely on the poor results coming from a recommender but rather the scientific metrics such as HR, NDCG, or MRR to calulate the health of the output after poisoning. Inside RAGPoison Lab each step is use as proof for the strength and weaknesses of the system before it gets to the model itself. each step of the workflow is an extra layer of risks and that is where the strength benchmark framework will be displayed

the issue that might be raised with this type of testing is whether the retrieval data can be evaluated without impacting the actual experiment. Each experiment is both ran on a clean and a poisoned index acting separately and never interacting with eachother giving us the opportunity to put probe alongevery step without inadvertently hurting our system. All experiment come from the same dataset source and then modified only for the attack. The results we will see later on come from the attack families we will use which won't affect the normal runs.

To make this possible RAGPoison Lab as said before sepparates the two concurrent runs in a confined setup through two different steams, one coming from the raw data, the other manipulated by our attacks. They are both hosted locally thanks to the Elasticsearch vector database. The framework also uses a very clear workflow on both streams: data preparation, attack creation, bulk index, experiment then log, metrics and summary outpput, all preserved in their JSON format \cite{harper2015movielens,ragpoisonlab_repo}.

Thanks to this implementation, the attack families are only used on one side of the experiement to avoid uncomperable or irreproducible runs. We are always capable of comparing with certitude the baseline to the poisonned because they both run at the same time on the same source. Since this framework is ran locally as a simulation, we can only discuss the results of this experiment as observations rather than a source of truth. It also avoids using real user data and risk tampering with real systems. The objective is to be transparent and not to try to prove that these results are general.

We wont try to discuss whether LLMs are dangerous in a context like this one, since it comes with it's benefits and inconviences. But we will discus how we can make an environnement where we can study the impact of this types of attacks through a system created for this exact reason. We will diiscuss the developpement and implementation hardships along side the proof coming from every step of our experimentation.

\subsection{Objectives}

The first objective of this thesis is to develop and use a controlled benchmark environement to understand the true security problem: what happens when the information retrieved by the recommender can no longer be trusted without getting into trying to prove that all LLM recommenders are unsafe. We will thus build a inspectable framework that will allow use  to compare the retrieval condition when the data is untouched and when it is affected without changing the surrounding environement. MovieLens 100K is prefect for this as it is small compared to other datasets but big enough to use in our scenario without involving live user or third party systems \cite{harper2015movielens}. RAGPoison Lab will be the framework that make this comparison possible while keeping it scope bound to what we want to achieve in this thesis, show all the process from preparation to output with metric probes and powerful and revealing logs \cite{ragpoisonlab_repo}.

The second objective is to try multiple attack families instead of relying on one type of attack. Targeted promotion asks whether we can manipulate a spc\ecific item to make it reach the recommended list. Untargeted degradation will answer whether general posioning of items will reduce the recommendation quality without trying to push a specific item. Prompt injection will raise an extra concern about instruction based payloads and if they impact the models ranking. The problem isnt only what items are retrieved, its also if the way the LLM innterprets certain modification will impact the output \cite{greshake2023indirectpi,yi2025bipia}. 

a Third objective is to mesure the recommendation failures through scientific metrics, for that we will use HR, NDCG, MRR and ASR. The first three will study the recommendation quality , while ASR will be only be meaningful for target oriented attacks \cite{herlocker2004evaluating,wang2013ndcg}. 

The fourth and last objective is to make sure this experiment is reproducible and ethical. We can already answer the ethical issue by mentionning the closed environnement behavior we selected with all tampering done locally, never to get accessed publicly. All actions are traced, all attacks are modular and configurable and each run gets its own configuration records and parameters saved along side it, allowing use to trust the scores we get, it will also help with the interpretation later on. This thesis will not attempt to solve poisoning or injections, but it will help us under stange each step in which a malicious actor could intervene and break the recommender system. Existing work on RAG and corpus poisoning display the importance of document level attack already, the objective here is to observe them on a deeper level in a simulated envrionement \cite{lewis2020rag,gao2024ragsurvey,zhong2023retrievalcorpora,zou2025poisonedrag_usenix}.

These objectives lead to four research questions. \textbf{RQ1:} How can a RAG-style recommendation pipeline be structured so clean and poisoned retrieval conditions remain comparable and auditable? \textbf{RQ2:} How do targeted promotion, prompt-injection-style poisoning, and untargeted degradation affect recommendation outputs in different ways? \textbf{RQ3:} Which combination of ranking metrics, attack-success metrics, traces, and configuration records is needed to interpret those effects without overstating them? \textbf{RQ4:} What practical defensive lessons become visible when the benchmark is studied as a full retrieval-to-ranking system rather than as a single model response?
